\section{The Secret Meaning of Marriage}

Recently several of Gustav Meyrink's books were republished: \emph{The White Dominican}, \emph{Walpurgis Night}, \emph{The Angel of the West Window}, after the publisher Bompiani had already published a new translation of the \emph{Golem}, a book that had in its time a great success in Germany and from which a film was also made.

\textbf{Gustav Meyrink} (1868-1932) is a truly unique author in his genre. The designation ``novel" is only approximately applied to his books. If the supernatural has a prominent part in them, the whole is not reduced, however, to the fantastic, as, for example, in Poe or Hofmann or even Lovecraft.

In them, the supernatural has, in its way, a realistic character through the frequent references to initiatic teachings, which sometimes are directly given, other times they come from symbolic contexts. For Meyrink this element is actually his primary purpose. So in an interview he declared that his novels were only ``coverings", that they had a symbolic content which reflected his experiences. He asserted, moreover, that he more than invented his characters and their events, he ``lived" them.

It is possible that Meyrink had an initiation in the proper sense. First of all, in his youth an impulse toward the supernatural had brought him to be interested in some spurious forms of ``occultism" and to the same spiritism from which, however, he turned away from them decisively, pronouncing severe judgments on them. Especially thanks to successive contacts with some Hindus and Cabalistic circles, he found the ``way" and was able to formulate a complex conception of life with a magical and initiatic orientation that, although expounded in novels—novels of a high artistic value — through his clarity and authenticity, it is difficult to find in works specifically dedicated to this material.

We can note here some of the principle aspects of this conception. As the center, we can value the doctrine of Awakening with the contrast between the habitual state of human existence and that of those who have passed onto a higher form of existence and who are also called the ``Living", in its eminent sense. If Buddhism at its origins itself had the doctrine of ``anatman", or the denial of a true I in common man, this is also Meyrink's view.

A character in Walpurgis Nights says:

\begin{quotex}
You truly believe that all those who often mill around in the streets possess an I. They do not even possess anything. They are rather possessed at every moment by a phantasm that plays the part of an I in them. 

\end{quotex}
Elsewhere Meyrink speaks of ``extinguished suns", of existence fundamentally spectral.

\begin{quotex}
Man is so firmly convinced of how much he is awake. In truth, however, he is imprisoned by a network of sleep and dreams that he himself has created. Those who are entangled pass through life like a herd led to the slaughterhouse. 

\end{quotex}
The man who breaks away from the herd, has found again that

\begin{quotex}
key of power over his lower nature, rusted since the flood, who is called to be awake. To be awake is everything.

Only the awakened man is, in reality, immortal. Stars and gods fade, he only remains and can accomplish everything that he wants. But there is a god above him. 

\end{quotex}
Concerning the way that leads to awakening, Meyrink speaks of the ``magical kingdom of thought", something different from the usual peculiar exercises of mental concentration pointed out by theosophy and ``spiritualists". In a text inserted in the Green Face, some techniques to lead to awakening are pointed out, techniques that in part recall those of Yoga and that start from the immobilization of the body in order thereby to detach from it and transcend a world of appearances and illusions. Elsewhere, he also mentions the test of remaining conscious in a change of state provoked by toxic fumes.

It seems, therefore, that Meryink also considers a type of predestination or vocation (in the sense of being called) necessary to reach toward Awakening. For example, he recalled the legends about the apparitions of beings who ``never died" — Elijah, John the Evangelist of the gnostics, the cabbalistic Chidher Green, to whom we could add the mysterious Islamic El Khidr — who would manifest themselves to those who ``were bitten by snakes". Or to lead to such is a ``destined" configuration of his own existence, a sudden interior reversal (that in a certain world brings to mind the ``satori" of Zen) or the sudden acceleration of existential rhythms: ``like a horse who up until now had gone by walking and in one stroke launches into a gallop".

The event, then, can also have tragic aspects, it can be that of a type of Walpurgis Night: forces that are liberated, that take possession of the being and transport him. The awakening happens after something like a nightmare, and on that very subtle line that separates the ``Way of Life" from the ``Way of Death".

In Meyrink's world, woman and the symbol of the androgyne have an important role. The words put into the mouth, in the Green Face, of the cabbalist Sephardi in speaking of the bridge that leads to Life are:

\begin{quotex}
No man can reach that goal alone; he needs a female companion. It is only possible, if at all, by a combination of male and female forces. Therein lies the secret meaning of marriage which had been lost to mankind for thousands of years. 

\end{quotex}
The reference is, however, especially to a particular ``magical" type of conjunction, and it is only to this that the symbol of the realization of the androgyne (of the complete being, more or less in the Platonic sense) is tied. The basic idea here is that the sexual instinct ``is the root of death", but that the task is not to eradicate it and to flee from woman as Christian asceticism proposes, but rather to absorb in man the feminine principle, separated on earth from the masculine principle, and a secret union that is not deprived of dangers.

The masculine and the feminine would not constitute just a polarity (which in common existence can also lead to a simple, banal complementarity). Between the two principles a tension and a latent but real antagonism exists. A character in the \emph{Angel of the West Window} contrasts the ``eros procreating like an animal" and common love, called ``plebeian", sexual relations in which the latent sexual polarity becomes manifest and extreme, as much to confer on experience a destructive character.

If in China woman was called ``the enemy" because she tends to capture the yang principle of man, Meyrink speaks of the ``draining death that comes from woman" in relation to the action that she would exert — insensibly and invisibly — already in a general way; with clear reference to tantric practices, he however considers — always in the clothing of fictional events — some procedures intended to make the ``magical element" carried by the masculine sexual energy not spill out and be lost in the feminine substance.

Among the various themes taken up by Meyrink from the initiatic Traditions, there is that of an ``Order", that may also correspond to the ``chain of the Living", i.e., of the ``Awakened Ones", and the theme of a supreme occult center of the world, the seat of beings who invisibly control the destinies of men. For the second theme, recurrent in one form or another in secret teachings and in the traditions of various cultures, we can return to the vast material collected and interpreted by \textbf{Rene Guenon} in his book \emph{The King of the World}.

The reader can also limit himself to the simple literary aspect of the ``novel" in Meyrink's works which, as we said, are unique in their genre.

Some readers will be able, however, to also value adequately the actual teachings, like those which we have brought attention to, woven into the narrative events and symbolized by them.

\begin{quotex}
This essay by \textbf{Julius Evola} was originally published in the journal ``Roma" on 7 July 1972 under the title \textit{Senso occulto del matrimonio}.
\end{quotex}


\flrightit{Posted on 2012-10-20 by Aeneas }

\begin{center}* * *\end{center}

\begin{footnotesize}\begin{sffamily}



\texttt{Audrey on 2012-10-20 at 20:07 said: }

Similar to Franz Bardon. Seems to be a beer putsh trend around the 19c. Meant not derogatorily, by the way.


\hfill

\texttt{Audrey on 2012-10-20 at 21:08 said: }

And he seem to touch on the same things as Franz Bardon, esp. in his the Key to the True Qabbalah. The Hindu influence remains a mystery to me at its roots, despite my acquaintance with it both in terms of

teaching as well as persons.


\hfill

\texttt{Audrey on 2012-10-21 at 01:18 said: }

Indeed! Hence for example, whereas I find rape absolutely unacceptable, as it robs the female of chances of growing conscious, it is true that in fact sexual instinct is te root of death. Evola wrote that even in death the man acts as a gateway for the female. Humorously, the power rape in the Army becomes subjugation to men's fear, cowardice, etc. Hence a perversity of survival and the entire food nonsense. Humorously, I carry an M16 which I shoot left-arm, even though I am right-handed. Nothing, and I mean nothing makes military males angrier, feeling weaker, etc. The whole notion that they cannot dominate a female's ass and fertilize it like it's mother earth drives them crazy. Sorry about the graphic quality of what is written. It is also why so many regard sex as initiatory and so forth. A female who has powers of life and death over men, and is not your mother/wife is extremely, how shall we say, unpalatable.


\hfill

\texttt{Tiffany Dawn Smith on 2021-10-20 at 17:14 said: }

What a great blog to have found today! I just ordered Guenon's King of the World on your recommendation. Thank you.


\end{sffamily}\end{footnotesize}
